%-------------------------
% Resume in LaTeX
% Modern sans-serif layout, ATS-friendly
%-------------------------

\documentclass[letterpaper,10pt]{article}

\usepackage{fontspec}
\setmainfont{Ubuntu}[
    Path=fonts/,
    Extension=.ttf,
    UprightFont=Ubuntu-Regular,
    BoldFont=Ubuntu-Bold,
    ItalicFont=Ubuntu-Italic,
    BoldItalicFont=Ubuntu-BoldItalic,
    Numbers=Monospaced
]
\usepackage{titlesec}
\usepackage{enumitem}
\usepackage[dvipsnames,table]{xcolor}
\usepackage[hidelinks]{hyperref}
\usepackage{fancyhdr}
\usepackage{tabularx}
\usepackage{fontawesome5}
\usepackage{geometry}
\usepackage{graphicx}
\usepackage{needspace}
\usepackage{microtype}

\emergencystretch=1.5em

\geometry{
    letterpaper,
    left=0.6in,
    right=0.6in,
    top=0.55in,
    bottom=0.55in
}

% Accent color
\definecolor{accent}{HTML}{1F4E79}
\definecolor{subtle}{HTML}{555555}
\definecolor{rule}{HTML}{BBBBBB}
\definecolor{logobg}{HTML}{ECECEC}
\definecolor{keyword}{HTML}{2F4F4F}

% Logo tile - light gray box so white-background logos stand out on the white page
\newcommand{\logobox}[1]{%
  {\setlength{\fboxsep}{2.5pt}%
   \colorbox{logobg}{\includegraphics[width=7mm,height=7mm,keepaspectratio]{#1}}}%
}

% Unbreakable bold keyword - no hyphenation inside tech terms
\newcommand{\tech}[1]{\textbf{\mbox{#1}}}

% Bold slate-gray keywords: \kw unbreakable (single token), \kwb breakable (phrase)
\newcommand{\kw}[1]{\textcolor{keyword}{\tech{#1}}}
\newcommand{\kwb}[1]{\textcolor{keyword}{\textbf{#1}}}

\hypersetup{
    colorlinks=true,
    urlcolor=accent,
    linkcolor=accent
}

\pagestyle{fancy}
\fancyhf{}
\renewcommand{\headrulewidth}{0pt}
\renewcommand{\footrulewidth}{0pt}
\fancyfoot[L]{\small\color{subtle}John Doe}
\fancyfoot[R]{\small\color{subtle}\thepage}

% Section formatting --- colored bold, thin rule
\titleformat{\section}
  {\Large\bfseries\color{accent}}
  {}{0em}
  {}
  [\vspace{2pt}\color{rule}\titlerule\color{black}]
\titlespacing*{\section}{0pt}{12pt}{10pt}

% Custom commands
% Indent width for body content (icon column)
\newlength{\bodyindent}
\setlength{\bodyindent}{12mm}

% #1=icon path  #2=company  #3=location  #4=role  #5=dates
\newcommand{\jobheading}[5]{%
  \vspace{8pt}
  \noindent
  \begin{minipage}[c]{\bodyindent}
    \centering
    \logobox{#1}
  \end{minipage}%
  \begin{minipage}[c]{\dimexpr\textwidth-\bodyindent}
    \textbf{\large #2}\hfill\textcolor{subtle}{#3}\\
    \textcolor{subtle}{#4}\hfill\textcolor{subtle}{#5}
  \end{minipage}
  \vspace{-4pt}
}

% Entry without icon, but aligned with iconned entries
% #1=title  #2=date
\newcommand{\plainheading}[2]{%
  \vspace{4pt}
  \noindent
  \hspace{\bodyindent}%
  \begin{minipage}[t]{\dimexpr\textwidth-\bodyindent}
    \vspace{0pt}
    \begin{tabular*}{\linewidth}{l@{\extracolsep{\fill}}r}
      #1 & \textcolor{subtle}{#2} \\
    \end{tabular*}
  \end{minipage}
  \vspace{-2pt}
}

% Entry with icon, single-line (for certs, projects)
% #1=icon  #2=title  #3=date
\newcommand{\iconentry}[3]{%
  \vspace{4pt}
  \noindent
  \begin{minipage}[c]{\bodyindent}
    \centering
    \logobox{#1}
  \end{minipage}%
  \begin{minipage}[c]{\dimexpr\textwidth-\bodyindent}
    #2\hfill\textcolor{subtle}{#3}
  \end{minipage}
  \vspace{-2pt}
}

\newcommand{\projheading}[2]{%
  \vspace{2pt}
  \begin{tabular*}{\textwidth}{l@{\extracolsep{\fill}}r}
    #1 & \textcolor{subtle}{#2} \\
  \end{tabular*}
  \vspace{-6pt}
}

% #1=icon  #2=school  #3=degree  #4=dates
\newcommand{\eduheading}[4]{%
  \vspace{4pt}
  \noindent
  \begin{minipage}[c]{\bodyindent}
    \centering
    \logobox{#1}
  \end{minipage}%
  \begin{minipage}[c]{\dimexpr\textwidth-\bodyindent}
    \textbf{#2}\hfill\textcolor{subtle}{#4}\\
    \textcolor{subtle}{#3}
  \end{minipage}
  \vspace{-2pt}
}

\setlist[itemize]{leftmargin=\dimexpr\bodyindent+1em, topsep=2pt, itemsep=1pt, parsep=0pt, label=\color{accent}\textbullet}

% Sub-role heading within a single job entry (e.g. dual-hat responsibilities)
\newcommand{\subrole}[1]{%
  \par\vspace{4pt}\noindent\hspace*{\bodyindent}%
  \textbf{\color{accent}\small #1}\par\vspace{-2pt}%
}

% Highlight tag row beneath job/project title --- appears outside bullet list
% #1=label (e.g. Stack, Acquired Skills)  #2=content
\newcommand{\stackline}[2]{%
  \par\vspace{3pt}
  \noindent\hspace*{\bodyindent}%
  \begin{minipage}[t]{\dimexpr\linewidth-\bodyindent}
    \small\textbf{\color{accent}#1:}~\textcolor{subtle}{#2}
  \end{minipage}\par
  \vspace{1pt}
}

\setlength{\parindent}{0pt}

\begin{document}

%-------- HEADER --------
\begin{center}
    {\huge \tech{\MakeUppercase{John Doe}}} \\[3pt]
    {\Large \color{accent} Staff Data Engineer} \\[6pt]
    {\small
        \faEnvelope\ \href{mailto:john.doe@example.com}{john.doe@example.com} \quad
        \faLinkedin\ \href{https://linkedin.com/in/johndoe}{/in/johndoe} \quad
        \faGithub\ \href{https://github.com/johndoe}{github.com/johndoe} \quad
        \faMapMarker*\ San Francisco, CA
    }
\end{center}

%-------- SUMMARY --------
\vspace{6pt}
{\small Staff Data Engineer with 8+ years designing and operating distributed data systems at scale. Led data platform initiatives serving 500+ engineers and petabyte-scale analytics workloads. Deep expertise in Apache Spark, Flink, Kafka, and real-time streaming architectures. Open-source contributor and maintainer of multiple data infrastructure projects.\par}

%-------- EXPERIENCE --------
\section{Experience}

\jobheading{icons/company.png}{Nexus Technologies}{San Francisco, CA}{Staff Data Engineer}{Jan 2022 -- Present}
\subrole{Data Platform --- Real-time analytics infrastructure}
\begin{itemize}
    \small
\item Designed and built a unified \kwb{real-time analytics platform} processing \kwb{2M+ events/sec} using Apache Kafka, Flink, and Pinot, serving sub-second OLAP queries across 12 product teams with 99.99\% uptime SLA.
    \item Migrated the company's \kwb{10 PB+ data lake} from Hive-managed Parquet on HDFS to \kwb{Apache Iceberg on S3} with incremental table maintenance, reducing storage costs by 40\% and improving query latencies through Iceberg's native partitioning and file pruning.
    \item Built an \kwb{internal data discovery and lineage platform} (Apache Atlas + custom enrichment layer) that indexed 15K+ datasets, tracked \kwb{column-level lineage} across 300+ pipelines, and reduced data incident MTTR from 4 hours to 15 minutes.
    \item Architected a \kwb{multi-region streaming replication layer} with Kafka MirrorMaker 2 and hybrid active-active conflict resolution, enabling the company to serve EU customers from ``Frankfurt with sub-50ms p99 latencies.
    \item Led the \kwb{data engineer interview loop redesign} — standardized rubric, system-design scenarios, and take-home evaluation — increasing hiring pipeline throughput by 60\% while maintaining quality bar.
\end{itemize}
\stackline{Stack}{Apache Kafka, Flink, Pinot, Iceberg, Spark, Airflow, S3, Terraform, Kubernetes, Python, Java}

\jobheading{icons/company.png}{StealthLayer}{New York, NY}{Senior Data Engineer}{Mar 2018 -- Dec 2021}
\subrole{Streaming Infrastructure}
\begin{itemize}
    \small
    \item Replaced a legacy batch ETL pipeline processing \kwb{500M events/day} with a \kwb{real-time streaming architecture} using Kafka Streams and Flink, reducing event-processing latency from 4 hours to under 30 seconds and enabling live fraud-detection models.
    \item Designed and deployed \kwb{Change Data Capture pipelines} (Debezium + Kafka Connect) ingesting from 40+ MySQL/Postgres shards, enabling near-real-time read replicas and removing the need for nightly batch dumps.
    \item Built a \kwb{cost-intelligence layer} that tagged every Spark job and S3 object with business-unit and team cost-centers; gave finance granular per-team billing and reduced cloud spend by 25\% through reclaimable-waste identification.
    \item Implemented \kwb{automated data quality framework} (Great Expectations + Airflow DAG orchestration) with 600+ expectations across 200 tables, cutting data-quality incidents from monthly to under 3 per quarter.
    \item Mentored 4 junior data engineers through structured onboarding, code review pairing, and internal tech talks; 2 were promoted within 18 months.
\end{itemize}
\stackline{Stack}{Kafka Streams, Flink, Debezium, Airflow, Great Expectations, Spark, S3, Terraform, Docker, Python, Scala}

\jobheading{icons/company.png}{RiverPoint Analytics}{Boston, MA}{Data Engineer}{Jun 2015 -- Feb 2018}
\begin{itemize}
    \small
    \item Built and maintained the company's \kwb{data warehouse and ELT pipeline} (Snowflake + dbt + Airflow) serving analytics and BI for 80+ internal stakeholders; designed dimensional models for finance, marketing, and product domains.
    \item Developed a \kwb{custom Python-based ETL framework} that abstracted away common ingestion patterns (API polling, S3 event-triggered loads, SFTP pulls) into reusable, testable plugins — reduced new pipeline development time by 70\%.
    \item Optimized \kwb{Spark batch jobs} processing 5 TB/night of clickstream data: tuned shuffle partitions, switched to columnar encoding, and implemented incremental processing — reduced run time from 6 hours to 45 minutes and cut cluster cost by 60\%.
    \item Created a \kwb{Snowflake cost-optimization dashboard} using Snowflake's ACCOUNT\_USAGE schema and Streamlit, giving engineering leads visibility into credit consumption by warehouse, query pattern, and team — drove a 30\% reduction in snowflake spend.
\end{itemize}
\stackline{Stack}{Snowflake, dbt, Airflow, Spark, Python, Streamlit, AWS (S3, EMR, Lambda), PostgreSQL, Redshift}

%-------- EDUCATION --------
\needspace{6\baselineskip}
\section{Education}
\eduheading{icons/edu.png}{Massachusetts Institute of Technology}{M.S., Computer Science --- Systems \& Data}{2013 -- 2015}
\hspace*{-2pt}\eduheading{icons/edu.png}{University of Illinois Urbana-Champaign}{B.S., Computer Engineering}{2009 -- 2013}
\vspace{4pt}

%-------- PROJECTS --------
\section{Projects}

\iconentry{icons/project.png}{\tech{Polaris} \textcolor{subtle}{$\cdot$ Declarative Flink streaming framework with auto-scaling} \quad \href{https://github.com/johndoe/polaris}{\faGithub\ GitHub}}{Jan 2025 -- Present}
\begin{itemize}
    \small
    \item Designed a \kwb{YAML-driven Flink DSL} (Parse $\rightarrow$ Optimize $\rightarrow$ Compile $\rightarrow$ Deploy) that lets teams define complex streaming topologies (joins, windowed aggregations, CDC pipelines) without writing a single line of Java or Scala.
    \item Built an \kwb{auto-scaling executor} that monitors Kafka consumer lag, Flink checkpoint duration, and backpressure metrics to dynamically adjust parallelism — reduced over-provisioning cost by 35\% with zero downtime during rescale.
    \item Implemented \kwb{pluggable state backends} (RocksDB, Redis, DuckDB) with transparent schema evolution, so users can migrate state stores without pipeline rebuilds.
    \item Shipped \kwb{native Prometheus metrics export} with pre-built Grafana dashboards for latency, throughput, watermark lag, and checkpoint health — adopted by 4 teams internally.
\end{itemize}
\stackline{Acquired Skills}{Apache Flink, Stream Processing, DSL Design, Kubernetes, Prometheus, Grafana, Open-Source}

\iconentry{icons/project.png}{\tech{Sensor} \textcolor{subtle}{$\cdot$ Cross-platform data quality \& lineage observability} \quad \href{https://github.com/johndoe/sensor}{\faGithub\ GitHub}}{Mar 2024 -- Dec 2024}
\begin{itemize}
    \small
    \item Designed a \kwb{pluggable quality rule engine} (expectation $\rightarrow$ validator $\rightarrow$ action) with 30+ built-in validators (nullness, uniqueness, distribution drift, freshness, referential integrity) expressed as version-controlled YAML policies.
    \item Built a \kwb{column-level lineage tracker} that intercepts Spark SQL plan nodes via Catalyst extensions, maps input-to-output column provenance, and emits OpenLineage-compatible events to Marquez — deployed across 80+ production pipelines.
    \item Implemented a \kwb{drifting metrics dashboard} (pre-computed statistical profiles stored in DuckDB, served via a FastAPI backend and a React frontend) that alerted teams when schema or distribution shifts exceeded configurable thresholds — cut data-incident detection time from 6 hours to under 5 minutes.
\end{itemize}
\stackline{Acquired Skills}{Spark SQL Catalyst, OpenLineage, Data Quality, React, FastAPI, DuckDB, Statistical Profiling}

%-------- CERTIFICATIONS --------
\needspace{5\baselineskip}
\section{Certifications}

\iconentry{icons/cert.png}{\textbf{Google Cloud Professional Data Engineer}}{Jan 2023}
\vspace{4pt}

\iconentry{icons/cert.png}{\textbf{AWS Certified Data Analytics --- Specialty}}{Mar 2021}
\vspace{4pt}

\iconentry{icons/cert.png}{\textbf{Databricks Certified Associate Data Engineer}}{Jun 2024}
\vspace{4pt}

%-------- SKILLS --------
\needspace{8\baselineskip}
\section{Skills}
\renewcommand{\arraystretch}{1.4}
\begin{tabularx}{\textwidth}{@{}>{\bfseries\color{accent}}l X@{}}
    Data Engineering    & Apache Spark, Flink, Kafka, Airflow, dbt, Iceberg, Delta Lake, Pinot, Trino \\
    Languages           & Python, Java, Scala, SQL, C++, Go, Bash \\
    Databases           & PostgreSQL, Snowflake, BigQuery, DuckDB, ClickHouse, Redis, MongoDB, Pinot \\
    Cloud \& DevOps     & AWS, GCP, Azure, Kubernetes, Docker, Terraform, Helm, GitHub Actions, Linux \\
    Streaming           & Kafka Streams, Flink, Debezium, KSQL, Pulsar, Kinesis \\
    Machine Learning    & Feature Stores, MLflow, Ray, XGBoost, scikit-learn, embedding pipelines \\
    Backend \& Tooling  & FastAPI, gRPC, GraphQL, Pydantic, Protobuf, Git, REST \\
    Observability       & Prometheus, Grafana, Datadog, OpenTelemetry, ELK, Jaeger \\
    Spoken Languages    & English (Native), Spanish (Professional), German (Conversational) \\
\end{tabularx}

\end{document}
